% !TEX root = ../macro.tex
\section{Real Business Cycle Model}\label{sec:rbc}

% BEGIN ORIGINAL RBC COURSE TEXT (pre-revision; comparison only)
% 
% Originated from Kyland and Prescott (1982), Long and Plosser (1983), RBC model shows great empirical importance. A standard RBC model\footnote{For a sound review over the topic of RBC, see Rebelo Sergio. 2005. “Real Business Cycle Models: Past, Present, and Future.” The Scandinavian Journal of Economics. } is in nature a Ramsey model with
% \begin{itemize}
% \tightlist
%   \item Endogenous (but constant in the long run, i.e. $n = 0$) labor supply
%   \item Stochastic productivity (i.e. real shocks)
% \end{itemize}
% 
% 
% \begin{blueblock}
% \begin{note}[Two principles of the RBC theory of business cycles]
% 
% 1. Money is of little importance in business cycles.
% 
% 2. Business cycles are created by rational agents responding optimally to
% real (not nominal) shocks - mostly fluctuations in productivity growth,
% but also fluctuations in government purchases, import prices, or preferences.
% \end{note}
% \end{blueblock}
% 
% 
% 
% 
% \sepline
% \subsection{UMP of household}
% A typical household entails product-consuming for $C_t$ and labor-providing for $N_t$. Only in this model, the agent value leisure $L_t=1-N_t$.
% 
% The UMP in RBC model is
% \begin{align}
% &\max\limits_{C_{t},N_{t},B_{t}} E_t \sum\limits_{t=0}^{\infty}\beta^{t}u(C_{t},N_{t})
% \\&
% s.t. \; P_{t}C_{t}+Q_{t}B_{t}\leqslant B_{t-1}+W_{t}N_{t}+D_{t}
% \text{, where }
% \begin{cases}
% \beta \in (0,1) \\
% Q_{t}=\frac{1}{1+i_{t}} \text{ is the price per bond bought today}
% \end{cases}
% \end{align}
% 
% Form a Lagrangian and acquire FOCs
% \begin{align}
% &\mathscr{L}=E_{0} \sum\limits_{t=0}^{\infty}\beta^{t}[u(C_{t},N_{t})+\lambda_{t}(P_{t}C_{t}+Q_{t}B_{t}-B_{t-1}-W_{t}N_{t}-D_{t} )]
% \\&
% \Rightarrow
% \begin{cases}\frac{\partial  \mathscr{L}}{\partial C_{t}}=0 \Rightarrow \beta^{t}[u_{C_{t}} -\lambda_{t}P_{t}]=0 \Rightarrow \frac{u_{C_{t}}}{E_{t}U_{c,t+1}}=\frac{\lambda_{t}}{E_{t}\lambda_{t+1}} \cdot \frac{P_{t}}{E_{t}{P_{t+1}}} \\
% \frac{\partial  \mathscr{L}}{\partial N_{t}}=0 \Rightarrow \beta^{t}[u_{N_{t}}+\lambda_{t} W_{t}] =0 \\
% \frac{\partial  \mathscr{L}}{\partial B_{t}}=0 \Rightarrow -\lambda_{t}Q_{t}\beta^{t}+E_{t}\lambda_{t+1}\beta^{t+1}=0 \Rightarrow \beta \frac{E_{t}\lambda_{t+1}}{\lambda_{t}}=Q_{t} \Rightarrow \frac{1}{Q_{t}}=1+i_{t}
% \end{cases}
% \end{align}
% 
% using these FOCs
% \begin{align}
% &\begin{cases}
%   \beta^{t}[u_{C_{t}} -\lambda_{t}P_{t}]=0 \\
%   \beta^{t}[u_{N_{t}}+\lambda_{t} W_{t}] =0
% \end{cases} \Rightarrow \color{red} \text{Labor supply equation} -\frac{u_{N_{t}}}{u_{C_{t}}}=\frac{W_{t}}{P_{t}}
% \\&
% \text{Also} \begin{cases}\frac{u_{C_{t}}}{E_{t}U_{c,t+1}}=\frac{\lambda_{t}}{E_{t}\lambda_{t+1}} \cdot \frac{P_{t}}{E_{t}{P_{t+1}}} \\
% \beta \frac{E_{t}\lambda_{t+1}}{\lambda_{t}}=Q_{t}
% \end{cases}
% \Rightarrow \color{red} \text{Euler's equation } \frac{Q_{t}}{\beta}=E_{t}\cdot
% \frac{\frac{U_{c,t+1}}{P_{t+1}}}{\frac{U_{c,t}}{P_{t}}}
% \end{align}
% 
% 
% \begin{greenblock}
% e.g. Let utility function be specific
% \begin{align}
% &
% u(C_{t},N_{t})=\begin{cases}
% \frac{c^{1-\sigma}_{t}-1}{1-\sigma}-\frac{N_{t}^{1+\varphi}}{1+\varphi},\sigma\neq 1 \\
% \log C_{t}-\frac{N_{t}^{1+\varphi}}{1+\varphi},\sigma=1\end{cases}
% \\&
% \Rightarrow \begin{cases}
% \frac{W_{t}}{P_{t}}=C_{t}^{\sigma}N_{t}^{\varphi} \\
% Q_{t}=\beta E_{t}\cdot (\frac{C_{t+1}}{C_{t}})^{-\sigma}\cdot \frac{P_{t}}{P_{t+1}}
% \end{cases}
% \end{align}
% \end{greenblock}
% 
% 
% 
% \sepline
% \subsection{PMP of firm}
% \begin{align}
% \min\limits_{N_{t}} W_t N_t \Leftrightarrow &\max\limits_{Y_t} P_{t}Y_{t}-W_{t}N_{t}
% \\&s.t. Y_t = A_t K_t^\alpha N_t^{1-\alpha}, \alpha \in [0,1)
% \\&
% \text{where }
% \begin{cases}
%   Y_{t}=A_{t}W_{t}^{1-\alpha}\\
%   A_{t}=A_{t-1}^{\rho\alpha} \epsilon_{t}^{\alpha}, \epsilon_{t}^{\alpha} \overset{iid}{\sim} N(0,\sigma_{\alpha}^{\epsilon})
%   \text{ (technology shocks)}
% \end{cases}
% \\&
% \Rightarrow \text{F.O.C.: } \frac{W_{t}}{P_{t}}=(1-\alpha)A_{t}N_{t}^{-\alpha} \text{ (labor demand equation)}
% \end{align}
% 
% \sepline
% \subsection{Economy described by RBC}
% {\color{red}
% The whole economy system is described as the balance of demand and supply
% \begin{align}
% &\text{AS}
% \begin{cases} Y_{t}=A_{t}N_{t}^{1-\alpha} \text{ (total production function)}, A_t=A^{\rho a}_{t-1} e^{\epsilon^a_t} \text{ (technology shocks)}\\
% \frac{W_{t}}{P_{t}}=C_{t}^{\sigma}N_{t}^{\varphi} \text{ (labor supply from UMP)}\\
% \frac{w_{t}}{P_{t}}=(1-\alpha)A_{t}N_{t}^{-\alpha} \text{ (labor demand from PMP)}
% \end{cases}
% \\&
% \text{AD}
% \begin{cases}Y_{t}=C_{t}+I_{t} \text{ (accouting equation)}\\
% Q_{t}=\beta E_{t} (\frac{C_{t+1}}{C_{t}})^{-a} \frac{P_{t}}{P_{t+1}}  \text{ (Euler equation from UMP)}\\
% \frac{M_{t}}{P_{t}}=\frac{Y_{t}}{Q_{t}^{-n}}  \text{ (demand of liquidity)}
% \end{cases}
% \end{align}
% }
% 
% 
% \sepline
% \subsection{Log-linearization}
% 
% Two methods we can transform.
% 1) The first is
% \begin{align}
% &
% \hat x_{t}=\ln x_{t}-\ln x^{*}
% \\&
% =\ln |\frac{ x_{t}}{x^{*}}|
% \\&
% =\ln |1+ \frac{x_{t}-x^{*}}{x^{*}}| \text{ (percentage of deviation from stable state)}
% \\&
% \approx \ln |1+\frac{1}{x^{*}}(x_{t}-x^{*})|
% \\&
% =\frac{x_{t}-x^{*}}{x^{*}}
% \end{align}
% 
% 2) Another approach would be
% \begin{align}
% &\ln x^{*}=ln x_{t}-\tilde{x_{t}}\Rightarrow \ln x_{t}=ln x^{*}+\tilde{x_{t}}
% \\&
% \Rightarrow e^{\ln x_{t}}=x_{t}=e^{n x^{*}+\tilde{x_{t}}}=x^{*}e^{\hat{x_{t}}} \approx x^{*} (1+\hat x_{t})
% \end{align}
% 
% Now we turn to the linearization of RBC model
% \begin{align}
% &
% \begin{cases}
% y_{t}=a_{t}+(1-\alpha)n_{t}
% \\
% y=a+(1-\alpha)n
% \end{cases}
% \Rightarrow \hat y_{t}=\hat a_{t}+(1-\alpha)\hat n_{t}
% \\&
% \Rightarrow
% \begin{cases}
% \hat w_{t}-\hat p_{t}=\hat a_{t}+(1-\alpha)\hat n_{t} \\
% \hat w_{t}-\hat p_{t}=a\hat c_{t}+\rho \hat n_{t}\end{cases}
% \Rightarrow \hat y_{t}=\hat c_{t}, \; m_{t}-\hat p_{t}=\hat y_{t} -\hat n_{it}
% \\&
% Y_{t}=C_{t}+I_{t}
% \\&
% Y\cdot e^{\hat y_{t}}=C\cdot e^{\hat c_{t}}+I\cdot e^{\hat i_{t}}
% \\&
% Y(1+\hat y_{t})=C(1+\hat c_{t})+I(1+\hat i_{t})
% \\&
% \hat y_{t}=\frac{C}{Y}\hat c_{t}+ \frac{I}{Y} \hat i_{t}\\&
% a_{t}=\rho_{t} a_{t-1}+\epsilon_{t}^{a}
% \\&
% 1=E_{t} e^{\ln \beta + +\ln \frac{1}{Q_{t}}-a(c_{t+1}-c_{t})-(P_{t+1}-P_{t})}
% \\&
% 1=E_{t} e^{\ln \beta + +\ln \frac{1}{Q_{t}}-a\triangle c_{t+1}-\pi_{t+1}}
% \\&
% 1=E_{t} I[1+(i_{t}-i)-a(\triangle c_{t+1}-\triangle c_{t}]-(\pi_{t+1}-\pi_{t})]
% \end{align}
% 
% {\color{red} Therefore a RBC model can be rewritten in linearized log equations
% \begin{align}
% \color{red}&\text{AS:} \begin{cases}
% \hat y_{t}=a_{t}+(1-\alpha)\hat n_{t} \text{ (output)}\\
% \hat w_{t}-\hat p_{t}=a_{t}-\alpha \hat n_{t} \text{ (labor demand)}\\
% \hat w_{t}-\hat p_{t}=a c_{t}+\varphi \hat n_{t} \text{ (labor supply)}\\
% \hat a_{t}=\rho_a \hat a_{t-1}+\epsilon_t^a \text{ (technology)}
% \end{cases}
% \\&
% \text{AD:} \begin{cases}
% \hat y_{t}=\hat c_{t} \text{ (product market clearing)}\\
% \hat c_{t}=E_{t}\hat c_{t+1}-\frac{1}{a}(\hat i_{t}-E_{t}\hat \pi_{t+1}) \\
% \hat m_{t}-\hat p_{t}=\hat y_{t}-\eta \cdot \hat i_{t} \text{ (money demand equation; $\eta$ is interest rate elasticity)}\\
% \hat r_{t}=i_{t}-E_{t} \hat \pi_{t+1} \text{ (Fisherian equation)}
% \end{cases}
% \end{align}
% }
% 
% \begin{blueblock}
% \begin{note}[Exceptation for Expectation]
% Notice this method does not apply to equations involving expectations.
% \end{note}
% \end{blueblock}
% 
% \begin{blueblock}
% \begin{note}[RBC has neutrality of money]
% The final answer of this standard RBC model is not relevant to any price variable, inferring the model being money-neutral.
% \end{note}
% \end{blueblock}
% 
% 
% 
% 
% 
% 
% 
% 
% 
% \subsection[Capital accumulation and investment (planned)]{Capital accumulation and investment}
% \textit{Planned section.} Extend the current fixed-capital benchmark to endogenous capital accumulation.
% Derive investment dynamics and distinguish the simplified benchmark from a
% full RBC model with a changing capital stock.
% END ORIGINAL RBC COURSE TEXT

\subsection{From a growth path to changing opportunities}

\subsubsection{What would make a Ramsey household change its working week?}

Imagine a group of factories. Their machines have not changed overnight, but
an interruption in power supply now makes each working shift less productive.
Orders, hours, investment, and consumption may all change. Start with questions
rather than labeling the event a ``technology shock.'' Can the factories produce
less with the same inputs? Are households poorer? Should they postpone work or
investment until conditions improve? Would the answer differ if the interruption
lasted one week rather than several years?

Ramsey already gives a way to compare consumption today and tomorrow. But with
fixed labor and a known production path, it leaves two margins closed. Households
cannot move working time toward more rewarding dates, and they never have to
choose before learning tomorrow's production conditions. Let labor be chosen;
let future productivity be uncertain. Keep preferences, ownership, production,
and market clearing explicit. This produces a tractable real business cycle
(RBC) benchmark.

This is a teaching reconstruction of the modeling questions, not a record of
any author's private thought process. Kydland and Prescott (1982) and Long and
Plosser (1983) are foundational contributions. Their actual environments differ;
the one-sector model below is a simpler teaching economy, without the former's
multi-period construction technology or the latter's sectoral structure.
\footnote{Finn E. Kydland and Edward C. Prescott (1982), ``Time to Build and
Aggregate Fluctuations,'' \textit{Econometrica}, 50(6), 1345--1370,
\url{https://www.minneapolisfed.org/~/media/files/research/prescott/papers/timetobuild.pdf};
John B. Long, Jr., and Charles I. Plosser (1983), ``Real Business Cycles,''
\textit{Journal of Political Economy}, 91(1), 39--69,
\url{https://doi.org/10.1086/261128}.}

\subsubsection{Choose the disturbance, then explain its transmission}

A scalar $A_t$ summarizes how much output given capital and hours can produce.
An unexpected change in it is an innovation. The persistence of $A_t$ describes
how long production opportunities change. Household choices and accumulated
capital determine how that change travels through the economy. These are three
different objects: innovation, persistence, and propagation.

A fall in $A_t$ need not mean that engineers forgot a blueprint. A power or
logistics disruption can reduce effective production opportunities. But treating
an event as a scalar productivity change is a modeling decision: energy inputs,
relative prices, network bottlenecks, and distributional losses have been
suppressed. A credit refusal or an absence of orders cannot automatically be
reclassified as technology. Their mechanisms need their own constraints.

The growth models explain accumulation and, in some cases, the source of trend
growth. Here we hold trend productivity constant to isolate fluctuations around
a stationary reference. This does not assert that actual output has no trend.
Adding a deterministic trend requires a consistent normalization of production,
preferences, and Euler conditions; it cannot be hidden inside the shock term.

\Needspace{11\baselineskip}
\begin{blueblock}
\textbf{A real benchmark is a conditional experiment.}
The original course principles become two questions. How much joint movement
can real disturbances produce when agents respond optimally? Which observed
movements require additional mechanisms? Money and nominal frictions are
excluded from the real allocation block below. Their exclusion does not
establish their empirical irrelevance. Government purchases, relative import
prices, and preference disturbances can be studied in extensions, but are not
silently included in this one-shock economy.
\end{blueblock}

\subsection{Complete conditions and the order of decisions}\label{sec:rbc-conditions}

The main model has endogenous capital. The course's fixed-capital benchmark
will be recovered as a separate restriction in Section~\ref{sec:rbc-fixed}.
Population is one and does not grow; hours are allowed to vary at every date.
In particular, zero population growth is not constant labor input.

\Needspace{18\baselineskip}
\begin{modelsetup}[1]
\begin{modelcondition}{01}{Dates, people, and the time endowment}
Time is discrete, $t=0,1,\ldots$. A closed economy contains many identical,
infinite-lived, price-taking households, aggregated to one representative
household of constant mass one. Each household has one unit of time per date:
$N_t\in[0,1]$ is work and $L_t=1-N_t$ is leisure. Work denotes an intensive
hours margin; there is no separate unemployment, participation, matching, or
involuntary rationing margin. Consumption satisfies $C_t>0$.
\end{modelcondition}
\begin{modelcondition}{02}{Goods, states, ownership, and timing}
One homogeneous good can be consumed or installed as capital at unit resource
cost. Households own the capital and firm shares. The given initial stock is
$K_0>0$; $K_t$ was chosen at $t-1$ and is usable at $t$. The current productivity
$A_t$ is observed before choosing $C_t,N_t,I_t,K_{t+1}$ and financial positions.
Next period's innovation is not yet observed. Capital is the only accumulated
productive stock; there are no inventories, skills, adjustment costs,
installation delays, or variable utilization in this benchmark.
\end{modelcondition}
\begin{modelcondition}{03}{Production and competitive firms}
Firms rent capital and hire labor in competitive spot markets. Their technology
is $Y_t=A_tK_t^\alpha N_t^{1-\alpha}$, with $A_t>0$ and $0<\alpha<1$.
It has constant returns in $(K_t,N_t)$, positive marginal products, and diminishing
marginal products holding the other input fixed. Firms take prices and $A_t$
as given; there is no markup, externality, or nominal price or wage rigidity.
\end{modelcondition}
\begin{modelcondition}{04}{The stochastic process and information}
Write $a_t=\log(A_t/\bar A)$, where $\bar A>0$, and let
\[
a_t=\rho_a a_{t-1}+\varepsilon_t^a,\qquad |\rho_a|<1,\qquad
E_{t-1}\varepsilon_t^a=0,\quad E[(\varepsilon_t^a)^2]=s_\varepsilon^2<\infty.
\]
Innovations are independent and identically distributed, with a distribution
known to households; Gaussian innovations are an optional specification.
$E_t$ conditions on the full observed history through $t$ and the equilibrium
rules, including current $K_t,A_t$. Rational expectations means using this
conditional law, not knowing the future realization. $a_t=0$ is the deterministic
reference, not a claim that $E A_t=\bar A$ under uncertainty.
\end{modelcondition}
\begin{modelcondition}{05}{Preferences and the household objective}
The household maximizes $E_0\sum_{t=0}^\infty\beta^t u(C_t,N_t)$, with
$0<\beta<1$, $\sigma>0$, $\chi>0$, and $\varphi>0$, where
\[
u(C,N)=
\begin{cases}
(C^{1-\sigma}-1)/(1-\sigma)-\chi N^{1+\varphi}/(1+\varphi),&\sigma\ne1,\\
\log C-\chi N^{1+\varphi}/(1+\varphi),&\sigma=1.
\end{cases}
\]
The original course normalization is $\chi=1$. $\sigma$ governs consumption
curvature; $\varphi$ governs the curvature of work disutility. We derive local
interior conditions with $0<N_t<1$. Parameter and shock choices must support
finite expected welfare and feasible paths; large shocks may activate corners
that a local solution does not describe.
\end{modelcondition}
\begin{modelcondition}{06}{Accumulation, prices, and household resources}
Capital depreciates at $0<\delta<1$ and satisfies
$K_{t+1}=(1-\delta)K_t+I_t\geq0$. The baseline permits reversible investment;
there is no additional $I_t\geq0$ constraint. Let $w_t$ be the real wage,
$q_t$ the real capital rental payment before depreciation, and $d_t$ real
profits distributed to households. With no financial trade the household budget is
\[
C_t+K_{t+1}=w_tN_t+(q_t+1-\delta)K_t+d_t.
\]
Households face no tax, transfer, borrowing wedge, or idiosyncratic uninsured
risk. Prices and returns are taken as given in individual optimization;
market clearing is imposed afterward.
\end{modelcondition}
\begin{modelcondition}{07}{Financial claims and nominal units}
For the course's nominal bond calculation, $P_t>0$ is the goods price,
$W_t=P_tw_t$, and $D_t=P_td_t$. $B_t$ pays one currency unit per unit at $t+1$
and costs $Q_t>0$ at $t$, so $R_t^n=1/Q_t=1+i_t$ is the gross nominal return.
The full budget is
\[
P_t[C_t+K_{t+1}-(1-\delta)K_t]+Q_tB_t
 =W_tN_t+P_tq_tK_t+B_{t-1}+D_t.
\]
Bonds are in zero aggregate net supply and $B_{-1}=0$. A household's small
financial trade is allowed when deriving its Euler condition; identical-agent
market clearing then sets aggregate $B_t=0$. A risk-free \emph{real} claim is
a separate hypothetical payoff, not the nominal bond deflated by a known
future price. There is no monetary asset in the core household budget.
\end{modelcondition}
\begin{modelcondition}{08}{Market clearing, solvency, and terminal behavior}
Household labor and capital supplies equal firms' demands. Goods clear as
$Y_t=C_t+I_t$, and constant returns with competitive pricing imply $d_t=0$.
Admissible household plans rule out debt rollovers that finance consumption
without resources. For the capital-only aggregate optimum the transversality
restriction is
\[
\lim_{T\to\infty}E_0[\beta^T C_T^{-\sigma}K_{T+1}]=0.
\]
Solvency restricts admissible borrowing; transversality selects optimal terminal
behavior. Neither follows merely from writing the interior Euler equations.
No-Ponzi restrictions also apply to any individual financial trades.
\end{modelcondition}
\begin{modelcondition}{09}{Equilibrium, reference state, and scope}
An equilibrium is a collection of history-contingent allocations and prices
satisfying household optimality and admissibility, firm optimality, the
productivity law, and all market-clearing conditions for the given initial
state. We select the locally stable interior solution near its deterministic
steady state. Money, government, foreign trade, monopoly, financial frictions,
and household heterogeneity are absent. A nominal policy or quantity closure
is an extra specification; the real equations alone do not select $P_t$.
\end{modelcondition}
\end{modelsetup}

The within-date sequence is
\[
\underbrace{K_t,\ A_t\text{ observed}}_{\text{inherited stock and current information}}
\ \longrightarrow\
\underbrace{N_t,\ Y_t,\ C_t,\ I_t}_{\text{work, production, and allocation}}
\ \longrightarrow\
\underbrace{K_{t+1}}_{\text{productive next date}}.
\]
An announcement about tomorrow can change today's choices; tomorrow's machine
stock cannot jump retroactively. In the experiment below only unanticipated
productivity innovations occur.

\subsection{Household choice: working now and transferring consumption}

\subsubsection{Retain the course's nominal bond route}

First suppress capital trade to isolate the original asset calculation. Write
nominal nonlabor income as $\mathcal D_t$; this includes the fixed factor's
income in the course benchmark. The budget is
\begin{equation}\label{eq:rbc-course-budget}
P_tC_t+Q_tB_t=B_{t-1}+W_tN_t+\mathcal D_t.
\end{equation}
Use resources minus expenditures in the multiplier term:
\begin{equation}
\mathscr L=E_0\sum_{t=0}^\infty\beta^t
\{u(C_t,N_t)+\lambda_t[B_{t-1}+W_tN_t+\mathcal D_t-P_tC_t-Q_tB_t]\}.
\end{equation}
For an interior choice at a history observed at $t$, the first two derivatives are
\begin{align}
u_{C,t}-\lambda_tP_t&=0,
 &\lambda_t&=u_{C,t}/P_t,\\
u_{N,t}+\lambda_tW_t&=0,
 &-u_{N,t}&=\lambda_tW_t.
\end{align}
The common factor $\beta^t$ and the probability of reaching this history cancel
from the two conditions. Divide the second condition by the first:
\begin{align}
\frac{-u_{N,t}}{u_{C,t}}&=\frac{\lambda_tW_t}{\lambda_tP_t},\\
-\frac{u_{N,t}}{u_{C,t}}&=\frac{W_t}{P_t}=w_t.\label{eq:rbc-labor-general}
\end{align}
The marginal disutility of another working hour equals its consumption-valued
wage. This is an intratemporal condition: all its objects are known at $t$.

One unit of $B_t$ costs $Q_t$ now and pays one currency unit in every next-date
state. Its two multiplier contributions give
\begin{align}
0&=-\beta^t\lambda_tQ_t+\beta^{t+1}E_t\lambda_{t+1},\\
\lambda_tQ_t&=\beta E_t\lambda_{t+1},\\
Q_t\frac{u_{C,t}}{P_t}&=\beta E_t\!\left[\frac{u_{C,t+1}}{P_{t+1}}\right],\\
Q_t&=\beta E_t\!\left[
\frac{u_{C,t+1}}{u_{C,t}}\frac{P_t}{P_{t+1}}\right].\label{eq:rbc-nominal-euler}
\end{align}
Current $P_t$ and $u_{C,t}$ can leave the expectation; future marginal utility
and the future price generally cannot. The last line retains the course's
nominal Euler equation with a consistent sign and information convention.

Differentiate the stated utility explicitly:
\begin{align}
\frac{d}{dC}\frac{C^{1-\sigma}-1}{1-\sigma}
 &=\frac{(1-\sigma)C^{-\sigma}}{1-\sigma}=C^{-\sigma},\\
\frac{d}{dC}\log C&=1/C=C^{-1},\\
\frac{\partial u}{\partial N}
 &=-\frac{\chi(1+\varphi)N^\varphi}{1+\varphi}=-\chi N^\varphi.
\end{align}
Substitution into the two optimality conditions yields
\begin{align}
-\frac{-\chi N_t^\varphi}{C_t^{-\sigma}}
 &=\chi C_t^\sigma N_t^\varphi=w_t,\label{eq:rbc-labor}\\
\frac{u_{C,t+1}}{u_{C,t}}
 &=\frac{C_{t+1}^{-\sigma}}{C_t^{-\sigma}}
   =\left(\frac{C_{t+1}}{C_t}\right)^{-\sigma},\\
Q_t&=\beta E_t\!\left[
\left(\frac{C_{t+1}}{C_t}\right)^{-\sigma}
\frac{P_t}{P_{t+1}}\right].\label{eq:rbc-bond-crra}
\end{align}
When $C_t$ and thus the marginal value of resources are held fixed,
\begin{align}
\log N_t&=\frac{\log w_t-\log\chi-\sigma\log C_t}{\varphi},\\
\left.\frac{\partial\log N_t}{\partial\log w_t}\right|_{u_{C,t}}
 &=1/\varphi.
\end{align}
This is the Frisch hours elasticity. Along an equilibrium, consumption also
changes: a higher wage creates both a reward to working and additional wealth.
The Frisch calculation alone does not determine the total hours response.

\Needspace{12\baselineskip}
\begin{blueblock}
\textbf{Keep random prices inside the expectation.}
In general $E_t(X/Z)\ne E_tX/E_tZ$, and $E_t(XZ)\ne E_tX\,E_tZ$.
For two equally likely future states with $(X,Z)=(1,1),(3,2)$,
\[
E(X/Z)=\tfrac12(1+3/2)=5/4,\qquad EX/EZ=2/(3/2)=4/3.
\]
The nominal bond pays a certain currency amount, but its purchasing power
$1/P_{t+1}$ is uncertain. A risk-free real bond instead has a certain goods
payoff. Confusing the two loses the connection between prices and consumption
risk.
\end{blueblock}

\subsubsection{Capital brings back the Ramsey margin, with uncertain returns}

For the full budget in card 07, consumption and hours have the same derivatives.
The cost of one extra unit of $K_{t+1}$ is $P_t$ at $t$. Next date it delivers
rental income $P_{t+1}q_{t+1}$ plus surviving capital $P_{t+1}(1-\delta)$.
Taking prices as given, the two-date derivative is
\begin{align}
0&=-\beta^t\lambda_tP_t
 +\beta^{t+1}E_t[\lambda_{t+1}P_{t+1}(q_{t+1}+1-\delta)],\\
\lambda_tP_t
 &=\beta E_t[\lambda_{t+1}P_{t+1}(q_{t+1}+1-\delta)],\\
C_t^{-\sigma}
 &=\beta E_t[C_{t+1}^{-\sigma}(q_{t+1}+1-\delta)].\label{eq:rbc-capital-euler}
\end{align}
Let $R^K_{t+1}=q_{t+1}+1-\delta$ and
$M_{t,t+1}=\beta(C_{t+1}/C_t)^{-\sigma}$. Division by current marginal utility gives
\begin{align}
1&=E_t[M_{t,t+1}R^K_{t+1}],\\
Q_t&=E_t[M_{t,t+1}P_t/P_{t+1}],\\
R_t^f&=1/E_tM_{t,t+1}.\label{eq:rbc-pricing}
\end{align}
The last line prices a hypothetical risk-free real payoff. The bond can help
price intertemporal transfers even though identical households hold zero bonds
in aggregate. It cannot make aggregate consumption exceed aggregate resources.

Expected capital returns need not equal the risk-free return. Expanding the
conditional covariance gives
\begin{align}
1&=(E_tM_{t,t+1})(E_tR^K_{t+1})
       +\operatorname{Cov}_t(M_{t,t+1},R^K_{t+1}),\\
E_tR^K_{t+1}
 &=R_t^f-\frac{\operatorname{Cov}_t(M_{t,t+1},R^K_{t+1})}{E_tM_{t,t+1}}.
\end{align}
A payoff arriving when consumption is scarce is especially valuable. This
risk distinction disappears to first order around the deterministic reference;
it is not absent from the exact model.

\subsection{Firms, resource clearing, and the real equilibrium}

\subsubsection{Differentiate before normalizing the capital stock}

In real goods the firm solves
\[
\max_{K_t^d,N_t^d}\;
 A_t(K_t^d)^\alpha(N_t^d)^{1-\alpha}-w_tN_t^d-q_tK_t^d.
\]
The household owns the stock; the firm rents its current services. Thus the
firm's capital demand is a within-date choice, even though aggregate capital
supply $K_t$ is predetermined. The input derivatives at an interior optimum are
\begin{align}
0&=(1-\alpha)A_t(K_t^d)^\alpha(N_t^d)^{-\alpha}-w_t,\\
0&=\alpha A_t(K_t^d)^{\alpha-1}(N_t^d)^{1-\alpha}-q_t.
\end{align}
Impose $K_t^d=K_t$ and $N_t^d=N_t$ afterward:
\begin{align}
w_t&=(1-\alpha)A_tK_t^\alpha N_t^{-\alpha}
       =(1-\alpha)Y_t/N_t,\label{eq:rbc-wage}\\
q_t&=\alpha A_tK_t^{\alpha-1}N_t^{1-\alpha}
       =\alpha Y_t/K_t.\label{eq:rbc-rent}
\end{align}
For a prescribed output $Y_t$, cost minimization over labor alone, with fixed
$K_t$, delivers the same conditional labor requirement. It is not an
unconstrained minimization of wage payments with no output requirement.
The stock cannot be replaced by the wage in the production function.

Factor income exhausts output:
\begin{align}
w_tN_t&=(1-\alpha)Y_t,\qquad q_tK_t=\alpha Y_t,\\
d_t&=Y_t-w_tN_t-q_tK_t
 =Y_t-(1-\alpha)Y_t-\alpha Y_t=0.
\end{align}
Zero profit does not remove households' capital income. Use the household
budget, zero bond supply, and these payments to obtain
\begin{align}
C_t+K_{t+1}&=w_tN_t+q_tK_t+(1-\delta)K_t+d_t,\\
 &=Y_t+(1-\delta)K_t,\\
Y_t-C_t&=K_{t+1}-(1-\delta)K_t=I_t,\\
Y_t&=C_t+I_t.\label{eq:rbc-resources}
\end{align}
Saving allocates current goods to capital; bond trade only rearranges claims
between households. The labor condition, resource constraint, and capital
Euler equation play different roles and must all hold.

\subsubsection{The planner route checks the same allocation}

With the stated competition, ownership, and absence of wedges, a planner can
choose the aggregate allocation under the same technology and information:
\begin{align}
\mathscr L^P=E_0\sum_{t=0}^\infty\beta^t\{u(C_t,N_t)
 +\mu_t[A_tK_t^\alpha N_t^{1-\alpha}+(1-\delta)K_t-C_t-K_{t+1}]\}.
\end{align}
Its derivatives are
\begin{align}
u_{C,t}-\mu_t&=0,\\
u_{N,t}+\mu_t(1-\alpha)A_tK_t^\alpha N_t^{-\alpha}&=0,\\
-\mu_t+\beta E_t[\mu_{t+1}
 (\alpha A_{t+1}K_{t+1}^{\alpha-1}N_{t+1}^{1-\alpha}+1-\delta)]&=0.
\end{align}
Substituting $\mu_t=C_t^{-\sigma}$ recovers
\eqref{eq:rbc-labor} and \eqref{eq:rbc-capital-euler}, with the competitive
factor prices. Concavity, feasibility, and the terminal restriction accompany
these interior conditions. The equivalence must be reconsidered when a wedge,
externality, markup, or market incompleteness changes the environment.

\Needspace{12\baselineskip}
\begin{greenblock}
\textbf{The equilibrium is a loop, not two independent curves.}
Observed productivity and inherited capital determine productive opportunities.
Households choose hours and consumption jointly, using both current wages and
expected future returns. Hours change output and factor prices; the consumption
choice leaves investment; investment changes tomorrow's capital. Prices must
support all of these choices at once. An AS/AD grouping can organize the
original equations, but cannot close the economy without its resource,
accumulation, expectations, and terminal conditions.
\end{greenblock}

A compact list of the \emph{exact} real conditions is therefore
\begin{align}
Y_t&=A_tK_t^\alpha N_t^{1-\alpha},\\
\chi C_t^\sigma N_t^\varphi&=(1-\alpha)Y_t/N_t,\\
C_t^{-\sigma}&=\beta E_t[C_{t+1}^{-\sigma}
 (\alpha Y_{t+1}/K_{t+1}+1-\delta)],\\
Y_t&=C_t+I_t,\\
K_{t+1}&=(1-\delta)K_t+I_t,\\
a_t&=\rho_a a_{t-1}+\varepsilon_t^a.
\end{align}
The initial state and transversality condition remain part of the model even
when they are not repeated in this list. Wages, rentals, and bond prices are
then recovered from their pricing equations. No nominal variable is required
to solve this real block.

\subsection{Recover the course's fixed-capital benchmark}\label{sec:rbc-fixed}

Freeze productive capital at $\bar K=1$, remove its investment choice, and
assume no depreciation, government purchases, or foreign trade. There is no
replacement expenditure, so $I_t=0$ and $C_t=Y_t$. This is a separate classroom
restriction, not the full model with its capital Euler equation still imposed.
Owners receive fixed-factor income $\mathcal D_t/P_t=\alpha Y_t$; the labor-only
firm's residual pays that factor. The course's $\chi=1$ may also be imposed.
The limit $\alpha=0$ recovers labor-only production, as allowed in the course;
the main capital model uses $0<\alpha<1$.

The real equations become
\begin{align}
Y_t&=A_tN_t^{1-\alpha},\qquad C_t=Y_t,\\
w_t&=(1-\alpha)A_tN_t^{-\alpha},\\
w_t&=\chi C_t^\sigma N_t^\varphi.
\end{align}
Solve this restricted allocation without skipping the wealth effect. Equate
wages and replace consumption by production:
\begin{align}
\chi C_t^\sigma N_t^\varphi
 &=(1-\alpha)A_tN_t^{-\alpha},\\
\chi(A_tN_t^{1-\alpha})^\sigma N_t^\varphi
 &=(1-\alpha)A_tN_t^{-\alpha},\\
N_t^{\sigma(1-\alpha)+\varphi+\alpha}
 &=\frac{1-\alpha}{\chi}A_t^{1-\sigma}.
\end{align}
Let $d_N=\sigma(1-\alpha)+\varphi+\alpha>0$. At an interior solution,
\begin{align}
N_t&=\left[\frac{1-\alpha}{\chi}A_t^{1-\sigma}\right]^{1/d_N},\\
\log N_t&=\frac{\log(1-\alpha)-\log\chi}{d_N}
              +\frac{1-\sigma}{d_N}\log A_t,\\
\frac{d\log N_t}{d\log A_t}&=\frac{1-\sigma}{d_N},\\
\frac{d\log Y_t}{d\log A_t}
 &=1+(1-\alpha)\frac{1-\sigma}{d_N}
 =\frac{1+\varphi}{d_N}.
\end{align}
The last equality follows by collecting
$d_N+(1-\alpha)(1-\sigma)=\varphi+\alpha+1-\alpha=1+\varphi$.
The formulas apply while $0<N_t<1$; they do not bypass the time endowment.

\Needspace{11\baselineskip}
\begin{redblock}
\textbf{Question: does better productivity necessarily raise work?}
In this fixed-factor economy, $\sigma=1$ gives
$N_t=[(1-\alpha)/\chi]^{1/(1+\varphi)}$, independent of $A_t$.
Higher productivity raises the wage and income, but the work-reward and wealth
margins exactly offset. For $\sigma<1$ hours rise; for $\sigma>1$ they fall,
provided the solution is interior. Endogenous labor is not by itself a
prediction that hours always move with output. The investment margin in the
full model changes this experiment.
\end{redblock}

There is no endogenous capital memory here. With serially independent $A_t$,
the real allocation depends only on the current shock. A persistent $A_t$
produces persistence mechanically; zero-net-supply bond prices do not create
another aggregate resource state. This is precisely why the planned capital
section matters.

\subsection{Capital accumulation and the deterministic reference}\label{sec:rbc-steady}

\subsubsection{Find the steady state in the order the conditions determine it}

Switch innovations off and set $A_t=\bar A$. A deterministic steady state has
constant $K^*,C^*,N^*,Y^*,I^*$. It is a reference for local expansion, not the
unconditional mean of the nonlinear stochastic economy. The Euler equation gives
\begin{align}
(C^*)^{-\sigma}&=\beta(C^*)^{-\sigma}(q^*+1-\delta),\\
1&=\beta(q^*+1-\delta),\\
q^*&=\beta^{-1}-1+\delta.
\end{align}
Let $v=Y^*/K^*=q^*/\alpha>0$. Capital stationarity and goods clearing imply
\begin{align}
K^*&=(1-\delta)K^*+I^*,&I^*&=\delta K^*,\\
C^*&=Y^*-\delta K^*,&s_I\equiv I^*/Y^*&=\delta/v,\\
s_C\equiv C^*/Y^*&=1-\delta/v,&s_C+s_I&=1.
\end{align}
Here $v>\delta$ since $q^*>\delta$ and $\alpha<1$, so $s_C>0$.
The steady capital--hours ratio $x=K^*/N^*$ follows from the rental equation:
\begin{align}
q^*&=\alpha\bar A(K^*)^{\alpha-1}(N^*)^{1-\alpha}
 =\alpha\bar A x^{\alpha-1},\\
x^{1-\alpha}&=\alpha\bar A/q^*,\\
x&=(\alpha\bar A/q^*)^{1/(1-\alpha)}.
\end{align}
Define $z=Y^*/N^*=\bar A x^\alpha$. Then $C^*=s_CzN^*$ and
$w^*=(1-\alpha)z$. The labor condition pins down the scale:
\begin{align}
\chi(s_CzN^*)^\sigma(N^*)^\varphi&=(1-\alpha)z,\\
(N^*)^{\sigma+\varphi}
 &=\frac{(1-\alpha)z}{\chi(s_Cz)^\sigma},\\
N^*&=\left[\frac{(1-\alpha)z}{\chi(s_Cz)^\sigma}\right]^{1/(\sigma+\varphi)},\\
K^*&=xN^*,\quad Y^*=zN^*,\quad C^*=s_CzN^*,\quad I^*=\delta xN^*.
\end{align}
If this $N^*$ is outside $(0,1)$, the interior formula is inapplicable. One can
instead choose an intended interior $N^*$ and solve the same equation for
$\chi$; that is a stated calibration, not an extra equilibrium condition.

\subsubsection{Why current investment and current capital must not be confused}

A factory can change operating hours after learning $A_t$, but the stock
$K_t$ is already installed. Investment uses goods today and changes
$K_{t+1}$. Even after $A_t$ returns to its reference, the economy can inherit
a different capital stock. A single innovation can therefore have effects
lasting beyond the disturbance itself. The equations must determine whether
the stock builds up or runs down; the accounting identity alone does not
choose investment.
